%!TEX root = ../thesis.tex
%*******************************************************************************
%****************************** Third Chapter **********************************
%*******************************************************************************
\chapter{EXPERIMENTAL AND RESULTS}

% **************************** Define Graphics Path **************************
\ifpdf
    \graphicspath{{Chapter4/Figs/Raster/}{Chapter4/Figs/PDF/}{Chapter4/Figs/}}
\else
    \graphicspath{{Chapter4/Figs/Vector/}{Chapter4/Figs/}}
\fi

\section{Datasets, Preprocessing, and Partitioning}
\subsection{Datasets and Data Preprocessing}
To evaluate the proposed framework across heterogeneous clinical distributions, we benchmark on three representative datasets from the MedMNIST v2 \cite{yang2023medmnist} standardized collection: PathMNIST, BloodMNIST, and OrganAMNIST. These benchmarks span four fundamental medical imaging modalities: histopathology slides, optical blood smear microscopy, volumetric abdominal computed tomography (CT), and surface dermatoscopy.

\begin{table}[htbp]
\centering
\caption{Summary of benchmark medical imaging datasets used in our federated evaluation.}
\label{tab:datasets}
\small
\begin{tabular}{llcrrrr}
\toprule
\textbf{Benchmark} & \textbf{Clinical Modality} & \textbf{Classes} & \textbf{Train} & \textbf{Val} & \textbf{Test} & \textbf{Total} \\
\midrule
PathMNIST    & Colon Histopathology  & 9  & 89,996 & 10,004 & 7,180  & 107,180 \\
BloodMNIST   & Blood Cell Microscopy & 8  & 11,959 &  1,712 & 3,421  &  17,092 \\
OrganAMNIST  & Abdominal CT          & 11 & 34,561 &  6,491 & 17,778 &  58,830 \\
\bottomrule
\end{tabular}
\end{table}

All images are standardized to a spatial resolution of $28 \times 28$ pixels. To ensure structural compatibility across models without modifying initial convolution layers, single-channel grayscale CT slices from OrganAMNIST are projected to three identical channels via channel replication ($\mathbf{x} \in \mathbb{R}^{3 \times 28 \times 28}$). Pixel intensities in $[0, 1]$ are linearly normalized to $[-1, 1]$ using channel-wise statistics $(\mu = 0.5, \sigma = 0.5)$.

\subsection{Federated Non-IID Partitioning}
Addressing severe data heterogeneity across healthcare institutions, we evaluate algorithmic resilience by partitioning datasets across $K = 12$ federated clients using two non-IID regimes: (i)Practical clinical partitioning and (ii) Dirichlet distribution partitioning.

\section{Baselines}
and here I write more \dots

\section{ Main Experiments and Results}
This section has been modified from ``Publication quality tables in \LaTeX*''
 by Simon Fear.

The layout of a table has been established over centuries of experience and 
should only be altered in extraordinary circumstances. 

When formatting a table, remember two simple guidelines at all times:

\begin{enumerate}
  \item Never, ever use vertical rules (lines).
  \item Never use double rules.
\end{enumerate}

These guidelines may seem extreme but I have
never found a good argument in favour of breaking them. For
example, if you feel that the information in the left half of
a table is so different from that on the right that it needs
to be separated by a vertical line, then you should use two
tables instead. Not everyone follows the second guideline:

There are three further guidelines worth mentioning here as they
are generally not known outside the circle of professional
typesetters and subeditors:

\begin{enumerate}\setcounter{enumi}{2}
  \item Put the units in the column heading (not in the body of
          the table).
  \item Always precede a decimal point by a digit; thus 0.1
      {\em not} just .1.
  \item Do not use `ditto' signs or any other such convention to
      repeat a previous value. In many circumstances a blank
      will serve just as well. If it won't, then repeat the value.
\end{enumerate}

A frequently seen mistake is to use `\textbackslash begin\{center\}' \dots `\textbackslash end\{center\}' inside a figure or table environment. This center environment can cause additional vertical space. If you want to avoid that just use `\textbackslash centering'


\begin{table}
\caption{A badly formatted table}
\centering
\label{table:bad_table}
\begin{tabular}{|l|c|c|c|c|}
\hline 
& \multicolumn{2}{c}{Species I} & \multicolumn{2}{c|}{Species II} \\ 
\hline
Dental measurement  & mean & SD  & mean & SD  \\ \hline 
\hline
I1MD & 6.23 & 0.91 & 5.2  & 0.7  \\
\hline 
I1LL & 7.48 & 0.56 & 8.7  & 0.71 \\
\hline 
I2MD & 3.99 & 0.63 & 4.22 & 0.54 \\
\hline 
I2LL & 6.81 & 0.02 & 6.66 & 0.01 \\
\hline 
CMD & 13.47 & 0.09 & 10.55 & 0.05 \\
\hline 
CBL & 11.88 & 0.05 & 13.11 & 0.04\\ 
\hline 
\end{tabular}
\end{table}

\begin{table}
\caption{A nice looking table}
\centering
\label{table:nice_table}
\begin{tabular}{l c c c c}
\hline 
\multirow{2}{*}{Dental measurement} & \multicolumn{2}{c}{Species I} & \multicolumn{2}{c}{Species II} \\ 
\cline{2-5}
  & mean & SD  & mean & SD  \\ 
\hline
I1MD & 6.23 & 0.91 & 5.2  & 0.7  \\

I1LL & 7.48 & 0.56 & 8.7  & 0.71 \\

I2MD & 3.99 & 0.63 & 4.22 & 0.54 \\

I2LL & 6.81 & 0.02 & 6.66 & 0.01 \\

CMD & 13.47 & 0.09 & 10.55 & 0.05 \\

CBL & 11.88 & 0.05 & 13.11 & 0.04\\ 
\hline 
\end{tabular}
\end{table}


\begin{table}
\caption{Even better looking table using booktabs}
\centering
\label{table:good_table}
\begin{tabular}{l c c c c}
\toprule
\multirow{2}{*}{Dental measurement} & \multicolumn{2}{c}{Species I} & \multicolumn{2}{c}{Species II} \\ 
\cmidrule{2-5}
  & mean & SD  & mean & SD  \\ 
\midrule
I1MD & 6.23 & 0.91 & 5.2  & 0.7  \\

I1LL & 7.48 & 0.56 & 8.7  & 0.71 \\

I2MD & 3.99 & 0.63 & 4.22 & 0.54 \\

I2LL & 6.81 & 0.02 & 6.66 & 0.01 \\

CMD & 13.47 & 0.09 & 10.55 & 0.05 \\

CBL & 11.88 & 0.05 & 13.11 & 0.04\\ 
\bottomrule
\end{tabular}
\end{table}
